% Analysis
% Target: 2 pages

\subsection{Feature Importance Analysis}
\label{subsec:feature_importance}

We conduct ablation studies to identify which feature groups contribute most to the 40.45\% improvement.

\subsubsection{Group-wise Ablation}

\begin{table}[t]
\centering
\small
\begin{tabular}{lccr}
\toprule
\textbf{Feature Set} & \textbf{Dim} & \textbf{Val Loss} & \textbf{Improvement} \\
\midrule
Text-Only (Baseline) & - & 0.1330 & - \\
\midrule
Full Acoustic & 46 & \textbf{0.0792} & \textbf{40.45\%} \\
\midrule
Prosodic only & 13 & 0.1233 & 7.29\% \\
Spectral only & 30 & 0.1228 & 7.67\% \\
\midrule
w/o Prosodic & 33 & 0.0988 & 25.71\% \\
\bottomrule
\end{tabular}
\caption{Feature group ablation on 1000 samples (3 epochs). Spectral features provide the primary contribution (25.71\%), while prosodic features enable synergy (+14.74\% additional gain).}
\label{tab:ablation}
\end{table}

\textbf{Key findings:}
\begin{enumerate}
    \item \textbf{Spectral features are the backbone}: w/o Prosodic (33D) achieves 25.71\% improvement, accounting for 63.5\% of total gain
    \item \textbf{Prosodic features enable synergy}: Adding prosodic (13D) to spectral (33D) yields 40.45\% total, a 14.74\% boost beyond additive expectation (7.29\% + 25.71\% = 33.00\%)
    \item \textbf{Synergy quantification}: 40.45\% - 33.00\% = \textbf{7.45\% synergy bonus} from prosodic-spectral interaction
\end{enumerate}

\subsubsection{Feature Synergy Analysis}

Table~\ref{tab:synergy} shows the non-linear interaction between feature groups.

\begin{table}[t]
\centering
\small
\begin{tabular}{lccc}
\toprule
\textbf{Combination} & \textbf{Expected} & \textbf{Actual} & \textbf{Synergy} \\
\midrule
Prosodic + Spectral & 33.00\% & 40.45\% & +7.45\% \\
\bottomrule
\end{tabular}
\caption{Feature synergy: actual performance exceeds additive expectation.}
\label{tab:synergy}
\end{table}

\textbf{Interpretation}: Prosodic features (pitch, energy contours) modulate spectral features (vocal tract resonance), creating temporal dynamics that enhance emotion encoding beyond what either feature set achieves alone.

\subsubsection{Individual Feature Analysis}
% TODO: Implement feature selection/SHAP analysis
\begin{itemize}
    \item Rank features by importance (mutual information, SHAP values)
    \item Identify minimal feature set achieving 95\% of full performance
    \item Analyze correlation between features
\end{itemize}

\subsection{Generalization Analysis}
\label{subsec:generalization}

\subsubsection{Cross-Domain Testing}
% TODO: Test on real human speech
\begin{itemize}
    \item \textbf{TTS $\to$ Real speech}: Test model trained on TTS on human recordings
    \item \textbf{Dataset}: IEMOCAP, RAVDESS emotion databases
    \item \textbf{Status}: Future work - requires human emotion speech datasets
\end{itemize}

\subsubsection{Speaker Variation}

We test cross-speaker generalization by training on one TTS speaker (Claribel Dervla, female) and testing on another (Damien Black, male). The test set contains 100 samples with identical text content but different speaker voice.

\begin{table}[t]
\centering
\small
\begin{tabular}{lcr}
\toprule
\textbf{Test Set} & \textbf{Val Loss} & \textbf{Degradation} \\
\midrule
In-domain (Claribel Dervla) & \textbf{0.0792} & - \\
Cross-speaker (Damien Black) & 3.0502 & +3752\% \\
\bottomrule
\end{tabular}
\caption{Cross-speaker generalization. Model shows catastrophic performance drop on different TTS speaker.}
\label{tab:cross_speaker}
\end{table}

\textbf{Critical Finding}: Acoustic features show \textbf{catastrophic generalization failure} across speakers (38.5$\times$ performance degradation). This reveals that the model learns speaker-specific voice characteristics rather than speaker-invariant emotion patterns.

\textbf{Root Cause Analysis}:
\begin{enumerate}
    \item \textbf{Absolute vs relative features}: Pitch features (mean, range) are speaker-dependent. Female voices (Claribel: $\sim$200-400 Hz) differ fundamentally from male voices (Damien: $\sim$100-150 Hz).
    \item \textbf{Missing normalization}: Features lack speaker normalization (z-score standardization within speaker).
    \item \textbf{Static aggregation}: Mean/std statistics lose temporal dynamics that convey emotion independently of voice timbre.
\end{enumerate}

\textbf{Implications}: This result challenges our hypothesis that hand-crafted features generalize better than learned representations. It suggests deep models like WavLM may learn more speaker-invariant patterns. Speaker normalization and relative features (pitch contours vs absolute values) are critical for generalization.

\subsection{Interpretability Analysis}
\label{subsec:interpretability}

\subsubsection{Feature Visualization}
% TODO: Create visualizations
\begin{itemize}
    \item t-SNE projection of acoustic features colored by emotion
    \item Compare with WavLM embeddings
    \item Show acoustic features form cleaner emotion clusters
\end{itemize}

\subsubsection{Case Study Analysis}
% TODO: Analyze specific examples
\begin{itemize}
    \item Select examples where acoustic > WavLM
    \item Analyze which acoustic features differ most
    \item Relate to speech science principles
\end{itemize}

\subsection{Computational Analysis}
\label{subsec:computational}

\begin{table}[t]
\centering
\small
\begin{tabular}{lrr}
\toprule
\textbf{Operation} & \textbf{WavLM} & \textbf{Acoustic} \\
\midrule
Feature extraction (ms/sample) & 145.2 & 8.3 \\
Memory footprint (MB) & 478 & 12 \\
Embedding size (KB/sample) & 1.0 & 0.18 \\
\bottomrule
\end{tabular}
\caption{Computational efficiency.}
\label{tab:computation}
\end{table}
